\documentclass[a4paper,12pt]{article}
\usepackage[utf8]{inputenc}
\usepackage{fullpage}
\usepackage[style=phys]{biblatex}
\usepackage[pdftex]{graphicx}
\usepackage{float}
\usepackage{hyperref}
\usepackage{cleveref}
\usepackage{physics}
\usepackage{siunitx}


\addbibresource{references.bib}

\newcommand{\efn}{e4$\nu$}

\newcommand{\Mu}{$\mu$}
\newcommand{\Tau}{$\tau$}

\newcommand{\Ne}{$\nu_e$}
\newcommand{\Nm}{$\nu_\mu$}
\newcommand{\Nt}{$\nu_\tau$}

\newcommand{\Wp}{$W^+$}
\newcommand{\Wm}{$W^-$}
\newcommand{\pip}{$\pi^+$}
\newcommand{\pim}{$\pi^-$}


\begin{document}
\pagestyle{empty}
\par\noindent\includegraphics[width=12cm]{PandA_crest.pdf}
\par\noindent

\vspace*{2cm}

\begin{center}
        \Large\bf \Large\bf Senior Honours Project \\
        \LARGE\bf Further analysis of \efn\ data
\end{center}

\vspace*{0.5cm}

\begin{center}
        \bf Jan Kocka \\
        February 2022
\end{center}
\vspace*{5mm}

\begin{abstract}
    In this project we continue the work on data from \efn.
    We further inspect the data generated with GENIE, mainly looking at the Delta resonances and compare then to data from the CLAS experiment.
\end{abstract}

\vspace*{1cm}

\subsubsection*{Declaration}

\begin{quotation}
  I declare that this project and report is my own work.
\end{quotation}

\vspace*{2cm}
Signature:\hspace*{8cm}Date:  31/10/2021

\vfill
{\bf Supervisor:} Dr. Cheryl Patrick
\hfill
10 Weeks
\newpage

\pagestyle{plain}
\setcounter{page}{1}

\tableofcontents

\newpage

\section{Background and Motivation}
This project concerns neutrinos and their properties, and so a brief description of the relevant properties and their scientific significance is presented.
Neutrinos are 3 uncharged leptons of different flavours (each corresponding to the other 3, charged, leptons -- electron, \Mu\ and \Tau\ of ascending masses).
Neutrinos interact only via the weak force and with very low cross sections, the Feynman diagrams of their interactions with the other fundamental particles through the W and Z bosons are shown in \cref{fig:nu_feyn}.

Their existence was originally proposed by Pauli in the 1930s to explain the beta decay, and their mass was long thought to be 0, most notably the standard model assumes this.
However at the end of the last century neutrinos have been observed to undergo oscillations which implies non-zero mass.
Since then they have been at the forefront of research, nowadays they are accepted to have mass and there are questions about them being their own antiparticles (Majorana) and possibly violating charge-parity symmetry.
The answers to these questions could significantly further our understanding of the universe and explain the dominance of matter over antimatter.

\begin{figure}[h]
    \centering
    \includegraphics{figures/NeutrinoFeynman.jpg}
    \caption{
        Feynman diagrams of weak force vertices for neutrino interactions, all other vertices including them can be obtain by suitable rotations of these
        (this requires replacing some particles with their antiparticles, including the \Wm becoming a \Wp).
        Say the first could be changed to represent a $e^-$ and \Wp\ colliding and forming a \Ne.
        Figures taken from \cite{potterFeynmanDiagramsParticlea}.
    }\label{fig:nu_feyn}
\end{figure}

\subsection{Neutrino oscillations}
In particle physics oscillations is a phenomenon that takes place when the particle mass eigenstates which govern its evolution in time are not the same as the state in which it is observed.
This phenomenon is not exclusive to neutrinos, the neutral Kaon is known to oscillate to its own antiparticle and vice-versa \cite{burkhardtWavelengthNeutrinoNeutral2003}.
% However it might seem a bit abstract, but as we know from basic quantum mechanics when we make an observation, the observed system will collapse into some particular state depending on the result of our observation.

In the case of neutrinos, we say there are 3 flavours and the way we differentiate them is by the weak interaction, so say an electron capture takes place when an electron interacts with a proton to form a neutron and a neutrino.
This is a weak interaction process and so far all interactions we have observed conserve the 3 separate lepton numbers (electron, \Mu\ and \Tau\ again) thus we observe that neutrino to be an electron neutrino and we can say it is collapsed into the electron neutrino eigenstate.

Further, that electron eigenstate can be expressed in terms of the mass eigenstates, and if the 2 eigenstate sets aren't identical it will have multiple non-zero components.
Then as the mass eigenstates travel differently (the neutrino has a definite $E$ and so through $E^2 = m^2 + p^2$ we get that the mass eigenstates have different momenta and so speeds), the relative components of the state at different positions as it travels will change periodically -- oscillate.
Thus the probabilities of measuring the neutrino in any of the 3 flavour states also oscillates.
% The result is that if a neutrino is measured to be of some known flavour the probability of measuring it at a later time in any of the other flavour states changes and it might later be measured to be of a different flavour.

A very important property resulting from the two eigenstate sets being different is the transformation matrix between them.
For neutrinos this is the PMNS\footnotemark\ mixing matrix, it is referred to as $U$ and is such that 
\begin{equation}
    \ket{\nu_\alpha} = \sum_i U_{\alpha i} \ket{\nu_i}, \text{ and }
    \ket{\nu_i} = \sum_\alpha U^\dag_{i \alpha} \ket{\nu_\alpha}
\end{equation}
where $\ket{\nu_\alpha}$ are the 3 flavour eigenstates and $\ket{\nu_i}$ the mass eigenstates.
% The exact form of this matrix then determines many of the neutrino properties, notably if all the matrix elements are real then they maintain Charge-Parity (CP) symmetry.
% \footnotetext{Stands for Pontecorvo–Maki–Nakagawa–Sakata matrix, it is the analogue of the CKM matrix for the mixing of quarks.}
Then to complete the picture, the probability of a neutrino measured to be of flavour $\alpha$ to be measured again in the flavour $\beta$ is given by 
\begin{equation}
    P(\alpha \rightarrow \beta) = \sum_i |U_{\alpha i} U^\dag_{i \beta}|^2 + 2\Re \sum_{j>i} U_{\alpha i} U^\dag_{i \beta} (U_{\alpha j} U^\dag_{j \beta})^* \exp(-i\frac{\Delta m^2_{ij}}{2}\frac{L}{E})
\end{equation}
where $E$ is it's energy, $L$ the distance between our measurements and $\Delta m^2_{ij}$ the mass difference of the mass eigenstates $i, j$ squared (for a more detailed analysis see \cite{zuberNeutrinoPhysics2020}).

\subsection{Experiments and Energy Reconstruction}
Finally, these probabilities are something we can measure, there are many active experiments working on this to determine the neutrino masses and PMNS matrix elements, the list includes Super-Kamiokande, OPERA, MINOS, DUNE, Hyper-Kamiokande.\footnote{**maybe change which ones I mention and add references to them}
In all of these experiments neutrinos from a source which has known fractions of the 3 flavours are beamed across large distances (about 1300\si{km} for DUNE, for oscillations of 1\si{GeV} \Nm\ $L$ of order $\sim 1000\si{km}$ is needed \cite{mezzettoThreeFlavorOscillationsAccelerator2020}) so that they oscillate and then some of the flavours (electron, \Mu\ or both) are measured.
This means we are measuring some of the probabilities above with a known $L$, however there aren't any suitable monoenergetic neutrino sources.
Experiments typically use either neutrinos from the sun, from the atmosphere due to the cosmic microwave background or from secondary neutrino beams at accelerators.
We naturally don't have any control over the solar and atmospheric neutrinos, and while we can to some degree control the energies of neutrinos from accelerator sources not to the necessary degree.

Because of that experiments rely on reconstructing incident neutrino energies.
As neutrinos are uncharged and of small mass they aren't detected directly, instead a large amount of some stable substance (argon, water or chlorine) is stored and then when a neutrino comes in it might interact with one of these atoms and produce detectable particles.
These always have to include a lepton of the corresponding flavour and often the

% Neutrino detectors rely on having a large amount of some substance (Argon, water, Chlorine) with with the neutrino then interacts.
% In many detectors (including water Cherenkov type) only the outgoing lepton

% We naturally don't have any control over the solar and atmospheric neutrinos, but for accelerator sources we can to some degree control the generated neutrino energies, but the control is limited and the neutrino flux is orders of magnitude lower \cite{zuberNeutrinoPhysics2020, deloguENUBETExperiment2022}.
% With NBBs we can know the neutrino energies to some level, however 


% All of these sources tend to have a wide energy range however, 
% However none of these sources have a dominant energy, they all have a wide range, but for calculations of mass differences and PMNS matrix elements we need the $E$ values.
% For that reasons experiments are forced to reconstruct the energies of incoming neutrinos based on the observed events.


% in all of these experiments the neutrinos are either from the sun, microwave cosmic radiation or secondary sources from accelerators\cite{zuberNeutrinoPhysics2020}.
% None of these 



\newpage


\section{Results \& Discussion}

\section{Conclusion}

\printbibliography

\end{document}
